\documentclass[main.tex]{subfiles}
\begin{document}

\color{black}

Mười tuần tiếp theo của stars*collection lại đưa chúng ta trở về với căn mancave của \hls, nơi những anh tài tiếp tục chứng minh rằng chỉ cần tập hợp họ lại với nhau, dù là chuyện nhỏ nhặt đến đâu cũng có thể trở thành một sự kiện náo loạn. Kaiho giờ đây đã quen thuộc hơn với vai trò tổng quản, nhưng điều đó không có nghĩa là cô sẽ được tận hưởng những ngày làm việc yên bình. Trái lại, những người đồng nghiệp nam của cô dường như ngày càng tìm được nhiều cách mới để thử thách sự kiên nhẫn của vị quản lý này.

Mọi chuyện bắt đầu với một cú mất điện bất ngờ tại căn mancave. Trong bóng tối, các anh tài chẳng những không chịu ngồi yên chờ điện được khôi phục mà còn thi nhau tạo ra những tình huống nhốn nháo, khiến nơi làm việc vốn đã chẳng mấy yên bình nay càng trở nên khó kiểm soát. Trong một diễn biến khác, Rikka lại gặp phải một sự cố có phần\dots\ khó nói khi phần đầu bị tách rời khỏi cơ thể, để mặc phần thân trở thành trò cười cho những người xung quanh. Chưa hết, Ryuuhou còn được hưởng một món quà sinh nhật vô cùng đặc biệt theo cách chẳng ai ngờ tới: một ca phẫu thuật mắt LASIK miễn phí, như món quà sinh nhật ``ngược'' mà Rikka dành tặng cô!

Những bí mật tưởng chừng có thể giấu kín cũng lần lượt bị phơi bày. Shien hết lần này đến lần khác phủ nhận chuyện bản thân mắc chứng viêm mũi dị ứng, đồng thời cư xử như thể cả căn mancave là nhà riêng của mình. Thế nhưng, màn tự biện hộ ấy chẳng thể kéo dài lâu khi Oga, Kaiho và Suzuno lần lượt đưa ra những bằng chứng khiến cậu khó lòng chối cãi. Còn khi thời tiết chuyển mùa, Miyabi lại phải nằm bẹp vì cảm lạnh. Các anh chàng nhân cơ hội ấy liên tục trêu ngươi cậu, nhưng cuối cùng lại phải trả giá khi chính cơn cảm lạnh cũng lần lượt quật ngã những người đã quá tự tin vào sức khỏe của mình.

Không chỉ vậy, những vấn đề cá nhân tưởng chừng chẳng đáng kể cũng có thể châm ngòi cho một cuộc khẩu chiến. Astel, vốn mang trong lòng nỗi niềm về chiều cao khiêm tốn của mình, đã có màn đối đáp không khoan nhượng với Roberu. Hai người cứ thế tranh luận qua lại, biến một chủ đề vốn chẳng có gì nghiêm trọng thành một cuộc đấu khẩu đầy tính giải trí. Trong khi đó, một trận chiến guitar lại mở ra theo cách không ai mong đợi: thay vì chỉ dùng đàn để chơi nhạc, các anh tài còn biến nó thành công cụ để\dots\ đập Aruran! Tình huống này thậm chí còn kéo theo cả một màn quảng cáo bảo hiểm cùng những câu chơi chữ khó đỡ, khiến người xem chẳng biết nên tập trung vào âm nhạc hay vào những gì đang diễn ra trước mắt.

Những món đồ kỳ quặc cũng tiếp tục xuất hiện trong căn mancave. Temma một ngày nọ vác về một hòn đá có kích thước khổng lồ, rồi còn rút ra hai thanh maracas vì tưởng chúng là kiếm. Màn nhầm lẫn này suýt nữa khiến cậu tái hiện lại cảnh Roberu từng bị một hòn đá rượt đuổi trong buổi giới thiệu chương trình. Cũng chính Temma cùng Miyabi lại có một ý tưởng khác thường không kém: tự chui vào thùng rác và tuyên bố muốn ``sống cuộc đời ốc sên". Chứng kiến hai người tự nguyện biến mình thành những cư dân trong thùng rác, Kaiho và Oga chỉ còn biết ngán ngẩm trước mức độ kỳ quặc ngày càng tăng của những người đồng nghiệp.

Nhưng có lẽ sự cố đáng nhớ nhất lại liên quan đến Rikka và một chút nước. Sau khi bị chập mạch, cậu bắt đầu kéo theo hàng loạt hiện tượng kỳ quái: những chiếc bánh bao mini được sản xuất hàng loạt từ cánh tay, một Rikka khác bất ngờ xuất hiện, và mọi thứ cứ thế vượt khỏi tầm kiểm soát. Tuy nhiên, tất cả hóa ra chỉ là giấc mơ mà Shien nhìn thấy sau khi lén ăn vụng chiếc bánh pudding của Rikka. Một cái kết vừa bất ngờ vừa hợp lý theo cách chẳng ai có thể giải thích, nhất là khi thủ phạm lại chính là người đã trải qua toàn bộ cơn ác mộng ấy!

Và nhắc đến những chuyện khó tin, không thể bỏ qua tài kể chuyện của Temma. Từ những sự cố có thật như bị trẹo cổ chân, bị Kaiho tẩm quất sau khi trẹo vai, cho tới cả chuyện trẹo cánh tay, cậu đều có thể biến chúng thành những câu chuyện dài với đủ loại tình tiết. Thế nhưng, ranh giới giữa sự thật và trí tưởng tượng trong lời kể của Temma đôi khi lại rất mong manh. Cậu thậm chí còn lôi cả Edison ra làm lý do để bao biện, rồi tiếp tục đưa người nghe đến một tiệm tóc có mức giá lên tới 85 nghìn yên. Chẳng biết nên gọi đó là tài kể chuyện, khả năng chém gió hay đơn giản chỉ là một thế giới riêng mà Temma tự tạo ra cho mình nữa!

Mười tuần tiếp theo lại trôi qua với những màn mất điện, những sự cố khó đỡ, những cuộc tranh cãi và vô số câu chuyện vượt xa trí tưởng tượng thông thường. Kaiho cũng đã có thêm thời gian để hiểu rõ hơn về tính cách của từng anh tài \hls, từ những người thích gây chuyện cho tới những người luôn vô tình trở thành tâm điểm của mọi trò đùa. Còn các chàng trai trong căn mancave cũng dần quen với sự hiện diện của vị tổng quản, dù điều đó không đồng nghĩa với việc họ sẽ bớt nghịch ngợm hay biết giữ trật tự hơn.

Hai quyển đầu tiên của stars*collection đã ghi lại những bước đi ban đầu của Kaiho trong vai trò tổng quản \hls, cùng quá trình cô dần hòa nhập với những con người có tính cách chẳng ai giống ai. Nhưng đây mới chỉ là khởi đầu. Phía trước vẫn còn nhiều tuần làm việc, nhiều câu chuyện mới và không ít tình huống mà có lẽ ngay cả Kaiho cũng chưa thể chuẩn bị tinh thần để đối mặt.

Hãy cùng đón chờ những câu chuyện tiếp theo trong **Quyển số Ba: Trò Lố Có Tổ Chức**, nơi Kaiho và các anh tài \hls sẽ\ tiếp tục đồng hành, tạo nên những kỷ niệm mới và tất nhiên, không thể thiếu những màn hỗn loạn mang đậm dấu ấn riêng của căn mancave!

\end{document}